\documentclass{article}
\usepackage[utf8]{inputenc}
\usepackage{a4wide}

\begin{document}

\section*{ร้านดังย่านบรรทัดทอง}

ผู้ปกครองของรัชชาพงษ์เปิดร้านอาหารอยู่ที่บรรทัดทองซึ่งขายดีมาก

ทำให้ต้องหาวิธีจัดการการรอชำระเงินเพื่อเพิ่มประสบการณ์การใช้งานแก่ลูกค้า

โดยรัชชาพงษ์ขอรับผิดชอบในส่วนการบริหารจัดการคิวลูกค้าในการเข้าไปจ่ายเงินกับคนคิดเงินด้วยตนเอง

ในกรณีที่มีลูกค้ารอชำระเงิน n คน

เมื่อลูกค้ามาเข้าคิวจะได้รับบัตรคิวโดยบัตรคิวแต่ละใบจะระบุตัวเลข 1 ถึง n

การให้บริการลูกค้าแต่ละคนจะไม่เท่ากัน ซึ่งต้องใช้เวลาตามที่กำหนดแทนด้วยค่า ai

และรัชชาพงษ์กำหนดให้คนคิดเงินแต่ละคนสามารถให้บริการลูกค้าได้สองคนพร้อมกัน

ซึ่งเวลาในการให้บริการทั้งสองคนคือเวลาที่มากที่สุดของสองคนที่ถูกเลือกให้เข้าบริการพร้อมกัน

เมื่อเข้ารับบริการชำระเงินพร้อมกันทั้ง 2 คนแล้ว จะเสร็จพร้อมกันเช่นเดียวกัน

(นั่นหมายความว่าอาจจะมีคนนึงที่ต้องรออีกคนเสร็จพร้อมกัน)



หากมีลูกค้ารอชำระเงินมากกว่าหนึ่งคน จะเลือกลูกค้าสองคนจากสามคนแรกมาบริการพร้อมกัน

แต่ถ้ามีลูกค้ารอชำระเงินแค่หนึ่งคน ก็ให้ลูกค้าคนนั้นไปที่เคาน์เตอร์

และใช้เวลาที่กำหนดในการให้บริการได้เลย ดังนั้นจำนวนรอบการให้บริการทั้งหมดจะเป็น ⌈n/2⌉

โดย n คือจำนวนลูกค้าที่รอชำระเงิน



\textbf{งานของคุณ:}ช่วยรัชชาพงษ์หาว่าต้องใช้เวลาน้อยที่สุดเท่าใดในการบริการลูกค้าที่รอชำระเงิน

n คน



ความต้องการเวลาในการให้บริการทั้งหมดจะเป็นเท่าไร

โดยให้บริการลูกค้าในแต่ละรอบตามลำดับที่เหมาะสมที่สุด

ซึ่งหมายถึงการเลือกคู่ลูกค้าที่จะใช้เวลาในการให้บริการน้อยที่สุดในแต่ละรอบ



\hfill\break



{\def\LTcaptype{none} % do not increment counter

\begin{longtable}[]{@{}

  >{\raggedright\arraybackslash}p{(\linewidth - 4\tabcolsep) * \real{0.3333}}

  >{\raggedright\arraybackslash}p{(\linewidth - 4\tabcolsep) * \real{0.3333}}

  >{\raggedright\arraybackslash}p{(\linewidth - 4\tabcolsep) * \real{0.3333}}@{}}

\toprule\noalign{}

\begin{minipage}[b]{\linewidth}\centering

input

\end{minipage} & \begin{minipage}[b]{\linewidth}\centering

output

\end{minipage} & \begin{minipage}[b]{\linewidth}\centering

comment

\end{minipage} \\

\midrule\noalign{}

\endhead

\bottomrule\noalign{}

\endlastfoot

\begin{minipage}[t]{\linewidth}\raggedright

4\\

1 2 3 4\strut

\end{minipage} & \begin{minipage}[t]{\linewidth}\raggedright

6\\

1 2\\

3 4\strut

\end{minipage} & \begin{minipage}[t]{\linewidth}\raggedright

คู่แรก: ลูกค้า 1 และ 2 บริการพร้อมกัน ใช้เวลา 2\\

คู่ที่สอง: ลูกค้า 3 และ 4 บริการพร้อมกัน ใช้เวลา 4 เวลาที่ใช้ทั้งหมดคือ 6\\

\strut

\end{minipage} \\

\begin{minipage}[t]{\linewidth}\raggedright

5\\

2 4 3 1 4\strut

\end{minipage} & \begin{minipage}[t]{\linewidth}\raggedright

8\\

1 3\\

2 5\\

4\strut

\end{minipage} & \begin{minipage}[t]{\linewidth}\raggedright

คู่แรก: ลูกค้า 1 และ 3 บริการพร้อมกัน ใช้เวลา 3\\

คู่ที่สอง: ลูกค้า 2 และ 5 บริการพร้อมกัน ใช้เวลา 4 ลูกค้า 4 ไปที่เคาน์เตอร์คนเดียว ใช้เวลา

1 เวลาที่ใช้ทั้งหมดคือ 8\\

\strut

\end{minipage} \\

\end{longtable}

}



\subsection*{Input}
บรรทัดที่ 1: จำนวนเต็ม 1 ตัว (n -- แทนจำนวนลูกค้าที่รอชำระเงิน) โดยที่ 1 ≤ n ≤ 1000



บรรทัดที่ 2: จำนวนเต็ม n ตัว a1, a2, ..., an (1 ≤ ai ≤ 106) แยกด้วยช่องว่าง

แทนเวลาที่ใช้ในการบริการลูกค้าแต่ละคนในการชำระเงิน โดยเรียงลำดับจากต้นคิวถึงท้ายของคิว



\subsection*{Output}
บรรทัดที่ 1: จำนวนเต็ม 1 ตัวแทน เวลาน้อยที่สุดเท่าใดในการบริการลูกค้าที่รอชำระเงิน n คน



บรรทัดที่ 2 จนถึง ⌈n/2⌉: แต่ละบรรทัดประกอบด้วยจำนวนเต็ม 2 ตัว แยกด้วยช่องว่าง

แสดงคิวของลูกค้าที่ได้เข้าไปรับบริการชำระเงิน ในกรณีที่ n เป็นจำนวนคี่

บรรทัดสุดท้ายจะมีแค่จำนวนเต็ม 1 ตัวเท่านั้น



\end{document}
